%%%%%%%%%% source scan idx 42 | folio 35 %%%%%%%%%%

2.31. 设 $C$ 是 $k$ 上算术亏格为 $p$ 的约化固有曲线。采用如下记号：

\begin{description}[leftmargin=2.7cm,style=nextline]
\item[$M_0$] $C$ 的形变的完备形式概形 $T$ 的维数；
\item[$M_1$] $\underline{\operatorname{Aut}}(C)$ 的 Lie 代数的维数。
\end{description}

对 $T$ 的一个不可约分支 $E$，用 $M_0^E$ 表示 $E$ 的维数。2.27 的整体对应结果如下。

\underline{命题} 2.32. \underline{若 $C$ 在 $E$ 的泛点上的形变光滑，则}

\begin{equation*}
3p-3=M_0^E-M_1.
\end{equation*}

设 $X\subset C(k)$ 是 $C$ 的光滑点的一个有限集合，使得 $C$ 上在每个点 $x\in X$ 消失的向量场必为零。$(C,X)$ 的形变函子可形式表出，且满足 2.28 的一个变体。$(C,X)$ 的形式模概形在 $T$ 上光滑，相对维数为 $\#X-M_1$，由 2.28 的类似结果得到

\begin{equation*}
M_0^E+(\#X-M_1)=3p-3+\#X,
\end{equation*}

这就证明了 2.32。

现在假定 $k$ 的\underline{特征为零}。则 $M_1=\dim\underline{\operatorname{Aut}}(C)$。下面给出 $m_1(s)$ 的类似解释。

%%%%%%%%%% source scan idx 43 | folio 36 %%%%%%%%%%

\underline{引理} 2.33. \underline{若 $k$ 的特征为零，则 $A_s$ 的每个导子均可延拓为 $\widetilde A_s$ 的导子。}

若 $D$ 是 $A_s$ 的导子，则

\begin{equation*}
\exp(Dt)=\sum_{n\geq0}\frac{D^nt^n}{n!}
\end{equation*}

是 $A_s\otimes_k k[t]$ 在 $(s,0)$ 处的完备化的自同构。通过结构的搬移，此自同构作用于该环的正规化 $(\widetilde A_s\otimes_k k[t])^{\wedge}$ 上。模 $t^2$ 约化，得到 $\widetilde A_s^{\wedge}\otimes_k k[t]/(t^2)$ 的自同构 $1+Dt$，从而 $D$ 延拓为 $\widetilde A_s^{\wedge}$ 的导子。由于它也可延拓到 $A_s$ 的分式域（当 $A_s$ 不整时则为全商环），所以它可延拓到 $\widetilde A_s$。

假设 $s$ 为奇点。类似的论证表明，由 $D$ 定义的 $\operatorname{Spec}(\widetilde A_s)$ 上的向量场在 $s$ 上方的每一点处消失。

设 $c$ 是 $\widetilde A_s$ 的理想，包含于 $\widetilde A_s/A_s$ 的导子理想中。置

\begin{equation*}
\widetilde G_c=\underline{\operatorname{Aut}}(\widetilde A_s/c),
\end{equation*}

并设 $G_c$ 是 $\widetilde G_c$ 中保持子代数 $A_s/c$ 不变的子群。若 $d\subset c$，则 $\widetilde G_d$ 保持 $c/d$ 不变且满射到 $\widetilde G_c$。注意，对于一致化参数 $t$，$g\in\widetilde G_c$ 由 $g(t)$ 决定，而 $g(t)$ 可以是任意一致化参数。因此，带基点齐性空间 $\widetilde G_c/G_c$ 与 $c$ 无关；将其记为 $\widetilde G/G$。
\begin{editorialnote}
上述关于任意 $d\subset c$ 的不变性断言需要限制。以结点 $A=k[[x,y]]/(xy)$ 为例，其正规化为 $k[[x]]\times k[[y]]$；取 $d=(x^3)\times(y^3)$、$c=(x)\times(y^2)$，则二者都包含于导子，交换两个分支给出模 $d$ 的自同构，却不保持 $c/d$。因此，原书这一论证不能直接用于任意这样的理想。若在完备正规化 $B=\prod_i k[[t_i]]$ 中选用各分支同次幂的理想 $c_N=\prod_i(t_i^N)$，并取 $N_0$ 使 $c_{N_0}$ 包含于导子，则 $N\geq N_0$ 的这些理想在导子内的有限余维理想中共尾。它们是根基的幂，受到全部自同构保持。对 $M\geq N\geq N_0$，截断自同构可以提升一致化参数的系数，故约化映射 $\pi_{MN}:\widetilde G_{c_M}\to\widetilde G_{c_N}$ 为满射。又因 $c_N\subset A$，$A/c_M$ 是 $A/c_N$ 的完整逆像，所以 $\pi_{MN}^{-1}(G_{c_N})=G_{c_M}$；因此对应的齐性商空间在这个受限的共尾系上典范同构。原法文第 36 页也未明写这一限制；此处保留其原论述并注明适用范围。
\end{editorialnote}

设 $\widetilde\tau_0$ 是 $\operatorname{Spec}(\widetilde A_s)$ 上在 $s$ 上方各点处消失的向量场构成的 $\widetilde A_s$-模。当 $s$ 为奇点时，有

\begin{equation*}
\dim_k(\widetilde\tau/\widetilde\tau_0)
=\text{分支数},
\end{equation*}

以及

\begin{equation*}
\widetilde\tau_0/\tau\overset\sim\longrightarrow
\text{商空间 }\widetilde G/G\text{ 在原点的切空间}.
\end{equation*}

%%%%%%%%%% source scan idx 44 | folio 37 %%%%%%%%%%

\underline{命题} 2.34. \underline{在特征零时，若 $s$ 为奇点，则整数 $m_1(s)$ 是分支数与 $\widetilde G/G$ 的维数之和。}

\bigskip
\begin{center}
\underline{\textsc{参考文献}}
\end{center}

\begin{description}[leftmargin=1.1cm,labelwidth=.8cm,labelsep=.3cm]
\item[{[1]}] P. Defrise, \textit{Étude locale des correspondances rationnelles
entre surfaces algébriques}，教授资格论文，Brussels，1949。

\item[{[2]}] P. Deligne 与 D. Mumford, ``The irreducibility of the space of
curves of a given genus,'' \textit{Publ. Math. IHES} \underline{36} (1969),
75--110.

\item[{[3]}] Enriques--Chisini, \textit{Lezioni sulla teoria geometrica delle
equazioni e delle funzioni algebriche}, Zanichelli, Bologna, 1924.

\item[{[4]}] R. Hartshorne, \textit{Residues and Duality}, Lecture Notes in
Mathematics \underline{20}, Springer, 1966.

\item[{[5]}] S. Lichtenbaum, ``Curves over discrete valuation rings,''
\textit{Amer. J. Math.} 90, no. 2（1968 年 4 月）, 380--405.

\item[{[6]}] I. Manin, ``Rational surfaces over perfect fields,''
\textit{Publ. Math. IHES}, no. 30, 55--114.

\item[{[7]}] D. Mumford, ``The topology of normal singularities of an
algebraic surface and a criterion for simplicity,'' \textit{Publ. Math. IHES},
no. 9, 5--22.

\item[{[8]}] A. Néron, \textit{Modèles minimaux des variétés abéliennes sur
les corps locaux et globaux}, Publ. Math., no. 21.

%%%%%%%%%% source scan idx 45 | folio 38 %%%%%%%%%%

\item[{[9]}] M. Raynaud, ``Spécialisation du foncteur de Picard. Critères
numériques de représentabilité,'' \textit{C. R. Acad. Sci. Paris} 264
（1967 年 6 月）, 1001--1004.

\item[{[10]}] Saint-Donat, \textit{Sur un théorème de G. Castelnuovo et
F. Enriques}，第三阶段博士论文，Lyon。

\item[{[11]}] O. Zariski, ``On polynomial ideals defined by infinitely near
base points,'' \textit{Amer. J. Math.} 60 (1938), 151.

\item[{[12]}] D. G. Northcott, ``On the genus formula for plane curves,''
\textit{J. London Math. Soc.} 30 (1955), 376--382.

\item[{[13]}] D. G. Northcott, ``Abstract dilatations and infinitely near
points,'' \textit{Proc. Cambridge Philos. Soc.} 52 (1956), 178--197.

\item[] 关于“无限邻近点”的概念，见下书第二章的 Lipman 附录。

\item[{[14]}] O. Zariski, \textit{Algebraic Surfaces}, Ergebnisse
\underline{61}, Springer-Verlag.
\end{description}
